% GENERATED FILE. Edit canonical lessons, then run npm run generate:books.
% canonical-source-hash: fnv1a64:6bf9a348fce2482c
% canonical-lessons: FR-C330-bouillir, FR-C330-ebullition, FR-C330-marmite, FR-C330-mijoter, FR-C330-cocotte

\chapter{Sur le feu}
\label{ch:fr-sur-le-feu}
\hlchaptermodality{330}

\begin{chapteropening}
\textbf{By the end of this chapter:} \emph{I can describe cooking on the hob: boiling, the boil, a big pot, simmering and the cast-iron pot.}

Complete the last lesson of chapter 330: I can describe cooking on the hob: boiling, the boil, a big pot, simmering and the cast-iron pot.
\end{chapteropening}

\section[bouillir]{bouillir — to boil}

\label{lesson:FR-C330-bouillir}



\begin{quote}
Before the new one: say the French for the whisk, then the French for to whisk.
\end{quote}

\subsection*{You'll want to know: bouillir}
\textbf{bouillir} — \textquotedblleft{}to boil\textquotedblright{}.

\begin{quote}
L'eau bout depuis cinq minutes.
\end{quote}

\begin{quote}
\emph{The water has been boiling for five minutes.}
\end{quote}

Irregular: \emph{l'eau bout}, \emph{elle a bouilli}. \emph{Faire bouillir} — to bring to the boil.

\begin{cousinweb}
\emph{bouillir} is inherited through Old French \emph{boillir} from Latin \emph{bullīre} \textquotedblleft{}to bubble, to boil\textquotedblright{} (EtymDB), built on \emph{bulla} \textquotedblleft{}a bubble\textquotedblright{}.
\end{cousinweb}

\begin{grammarlens}[title={What you've built}]
One more word for telling, explaining and describing at B1.
\end{grammarlens}

\subsection*{Your turn}
\begin{itemize}
\raggedright
  \item \emph{Say it:} \emph{bouillir}
  \item \emph{Say it:} \emph{bouillir}, once more
  \item \emph{Say it:} a sentence of your own with \emph{bouillir}
  \item \emph{Recall:} say the French for the whisk, then the French for to whisk, then say \emph{bouillir} again
\end{itemize}

\begin{tcolorbox}[breakable,colback=teal!4,colframe=teal!35!black,title={Before you move on}]
What does bouillir mean? (\textquotedblleft{}To boil\textquotedblright{}.) Say it once more.
\end{tcolorbox}
\section[l'ébullition]{l'ébullition — boiling, the boiling point}

\label{lesson:FR-C330-ebullition}



\begin{quote}
Before the new one: say the French for to whisk, then the French for to boil.
\end{quote}

\subsection*{You'll want to know: l'ébullition}
\textbf{l'ébullition} — \textquotedblleft{}boiling, the boiling point\textquotedblright{}.

\begin{quote}
Il met les pâtes quand l'eau arrive à ébullition.
\end{quote}

\begin{quote}
\emph{He puts the pasta in when the water comes to the boil.}
\end{quote}

\emph{Porter à ébullition} — to bring to the boil.

\begin{cousinweb}
\emph{ébullition} is a learned borrowing from Late Latin \emph{ēbullītiō} \textquotedblleft{}a bubbling up\textquotedblright{} (EtymDB), built on \emph{bulla} \textquotedblleft{}a bubble\textquotedblright{}.
\end{cousinweb}

\begin{grammarlens}[title={What you've built}]
One more word for telling, explaining and describing at B1.
\end{grammarlens}

\subsection*{Your turn}
\begin{itemize}
\raggedright
  \item \emph{Say it:} \emph{l'ébullition}
  \item \emph{Say it:} \emph{l'ébullition}, once more
  \item \emph{Say it:} a sentence of your own with \emph{l'ébullition}
  \item \emph{Recall:} say the French for to whisk, then the French for to boil, then say \emph{l'ébullition} again
\end{itemize}

\begin{tcolorbox}[breakable,colback=teal!4,colframe=teal!35!black,title={Before you move on}]
What does l'ébullition mean? (\textquotedblleft{}Boiling, the boiling point\textquotedblright{}.) Say it once more.
\end{tcolorbox}
\section[la marmite]{la marmite — the cooking pot, the stockpot}

\label{lesson:FR-C330-marmite}



\begin{quote}
Before the new one: say the French for to boil, then the French for boiling.
\end{quote}

\subsection*{You'll want to know: la marmite}
\textbf{la marmite} — \textquotedblleft{}the cooking pot, the stockpot\textquotedblright{}.

\begin{quote}
La soupe cuit dans une grande marmite.
\end{quote}

\begin{quote}
\emph{The soup is cooking in a big pot.}
\end{quote}

A large pot with a lid and two handles.

\begin{cousinweb}
The origin of \emph{marmite} is not settled; an Old French adjective \emph{marmite} \textquotedblleft{}hypocritical\textquotedblright{} is listed beside it (EtymDB), perhaps a pot that hides what it holds.
\end{cousinweb}

\begin{grammarlens}[title={What you've built}]
One more word for telling, explaining and describing at B1.
\end{grammarlens}

\subsection*{Your turn}
\begin{itemize}
\raggedright
  \item \emph{Say it:} \emph{la marmite}
  \item \emph{Say it:} \emph{la marmite}, once more
  \item \emph{Say it:} a sentence of your own with \emph{la marmite}
  \item \emph{Recall:} say the French for to boil, then the French for boiling, then say \emph{la marmite} again
\end{itemize}

\begin{tcolorbox}[breakable,colback=teal!4,colframe=teal!35!black,title={Before you move on}]
What does la marmite mean? (\textquotedblleft{}The cooking pot, the stockpot\textquotedblright{}.) Say it once more.
\end{tcolorbox}
\section[mijoter]{mijoter — to simmer; to cook up (a plan)}

\label{lesson:FR-C330-mijoter}



\begin{quote}
Before the new one: say the French for boiling, then the French for the cooking pot.
\end{quote}

\subsection*{You'll want to know: mijoter}
\textbf{mijoter} — \textquotedblleft{}to simmer; to cook up (a plan)\textquotedblright{}.

\begin{quote}
La soupe mijote doucement.
\end{quote}

\begin{quote}
\emph{The soup is simmering gently.}
\end{quote}

Informal: \emph{qu'est-ce que tu mijotes ?} — what are you up to?

\begin{cousinweb}
The origin of \emph{mijoter} is not settled; it is usually linked to an old dialect word for a place where fruit ripens (EtymDB lists no parent).
\end{cousinweb}

\begin{grammarlens}[title={What you've built}]
One more word for telling, explaining and describing at B1.
\end{grammarlens}

\subsection*{Your turn}
\begin{itemize}
\raggedright
  \item \emph{Say it:} \emph{mijoter}
  \item \emph{Say it:} \emph{mijoter}, once more
  \item \emph{Say it:} a sentence of your own with \emph{mijoter}
  \item \emph{Recall:} say the French for boiling, then the French for the cooking pot, then say \emph{mijoter} again
\end{itemize}

\begin{tcolorbox}[breakable,colback=teal!4,colframe=teal!35!black,title={Before you move on}]
What does mijoter mean? (\textquotedblleft{}To simmer; to cook up (a plan)\textquotedblright{}.) Say it once more.
\end{tcolorbox}
\section[la cocotte]{la cocotte — the casserole dish, the cast-iron pot}

\label{lesson:FR-C330-cocotte}



\begin{quote}
Before the new one: say the French for the cooking pot, then the French for to simmer.
\end{quote}

\subsection*{You'll want to know: la cocotte}
\textbf{la cocotte} — \textquotedblleft{}the casserole dish, the cast-iron pot\textquotedblright{}.

\begin{quote}
Elle fait cuire le poulet dans une cocotte.
\end{quote}

\begin{quote}
\emph{She cooks the chicken in a casserole dish.}
\end{quote}

\emph{La cocotte-minute} — the pressure cooker. In baby talk, \emph{une cocotte} is a hen.

\begin{cousinweb}
The origin of \emph{cocotte} is not settled: one sense is a child's word for a hen, the other an older word for a pot (EtymDB lists no parent).
\end{cousinweb}

\begin{grammarlens}[title={What you've built}]
One more word for telling, explaining and describing at B1.
\end{grammarlens}

\subsection*{Your turn}
\begin{itemize}
\raggedright
  \item \emph{Say it:} \emph{la cocotte}
  \item \emph{Say it:} \emph{la cocotte}, once more
  \item \emph{Say it:} a sentence of your own with \emph{la cocotte}
  \item \emph{Recall:} say the French for the cooking pot, then the French for to simmer, then say \emph{la cocotte} again
\end{itemize}

\begin{tcolorbox}[breakable,colback=teal!4,colframe=teal!35!black,title={Before you move on}]
What does la cocotte mean? (\textquotedblleft{}The casserole dish, the cast-iron pot\textquotedblright{}.) Say it once more.
\end{tcolorbox}
